\chapter{Conclusion et perspectives}
\label{chap:conclusion}

\section{Bilan}

Le projet a livré un aquarium animé fonctionnel sous Linux et
GTK4, peuplé par défaut de huit espèces et offrant des
interactions prédateur/proie crédibles. Les objectifs initiaux
sont atteints~:

\begin{itemize}
  \item Le moteur de simulation à pas fixe est stable et
    déterministe, avec les trois forces de Reynolds
    correctement implémentées et complétées par des
    comportements de fuite, de chasse et d'évitement
    intra-espèce.

  \item Le cycle de vie complet des entités (\code{ALIVE},
    \code{DYING}, \code{INACTIVE}) est en place, avec un
    recyclage qui borne statiquement la consommation mémoire.

  \item Le \emph{système de configuration externe} en TOML
    permet effectivement à un utilisateur non développeur
    d'ajuster la simulation, d'ajouter ou de retirer des
    espèces, sans recompilation. Les chemins de
    \emph{fallback}, le rechargement à chaud et la gestion
    d'erreurs lisible répondent aux attentes posées au
    chapitre~\ref{chap:analyse}.

  \item L'interface graphique propose les contrôles essentiels~:
    pause, vitesse globale, sélection du \emph{pool}, ajout et
    retrait précis de poissons, plein écran, et rechargement de
    la configuration.
\end{itemize}

Le tableau~\ref{tab:bilan} synthétise les exigences
fonctionnelles et leur statut.

\begin{table}[h!]
\centering
\footnotesize
\begin{tabular}{@{}llc@{}}
\toprule
\textbf{Réf.} & \textbf{Exigence} & \textbf{Statut} \\
\midrule
EF1  & Simulation de banc multi-espèces       & livré \\
EF2  & Prédateur/proie avec mouth-based eating & livré \\
EF3  & Fuite des proies + adrénaline          & livré \\
EF4  & Cycle de vie 3 états                    & livré \\
EF5  & Famine                                  & livré \\
EF6  & Configuration TOML externe              & livré \\
EF7  & Recherche multi-niveaux                 & livré \\
EF8  & Ajout/suppression sans recompilation    & livré \\
EF9  & Validation et messages d'erreur         & livré \\
EF10 & Rechargement à chaud                    & livré \\
EF11 & Contrôles utilisateur                   & livré \\
EF12 & Plein écran                             & livré \\
\bottomrule
\end{tabular}
\caption{Statut des exigences fonctionnelles.}
\label{tab:bilan}
\end{table}

\section{Limites techniques}

Le projet, dans son état actuel, présente plusieurs limites
identifiées~:

\paragraph{Complexité quadratique.} Toutes les détections de
voisinage parcourent l'ensemble des entités, ce qui rend le coût
$O(N^2)$ par tick. Au-delà de quelques centaines d'entités, cela
deviendrait limitant. Une structure spatiale (grille uniforme,
\emph{spatial hashing}) ramènerait le coût à $O(N \log N)$ voire
$O(N)$ dans le cas amorti.

\paragraph{Rendu widget par widget.} Chaque entité est un widget
GTK distinct, soumis à un cycle complet de \emph{layout} et de
\emph{styling} CSS à chaque tick. Cela limite le nombre maximal
d'entités tenable à 60~Hz. Un rendu via \code{GtkSnapshot} ou un
\code{GtkDrawingArea} unique avec \code{cairo} serait beaucoup
plus efficace.

\paragraph{Validation post-chargement minimale.} Le parseur
TOML ne vérifie pas la cohérence sémantique des paramètres~: un
\code{max\_speed} à zéro, un \code{flock\_size} négatif ou un
\code{respawn\_delay} infini ne déclencheraient pas d'erreur. Une
passe de validation explicite avec plages tolérées serait utile.

\paragraph{Pas de tests automatisés.} Les régressions sont
détectées par observation visuelle. Un harnais sur les fonctions
pures (\code{vec2\_*}, parsing TOML, résolution \code{prey\_mask})
serait peu coûteux à mettre en place.

\paragraph{Pas d'instrumentation de performance.} Aucun
compteur FPS, aucune mesure du coût d'un tick.

\section{Perspectives}

Plusieurs pistes d'évolution se dégagent naturellement~:

\paragraph{Compteur FPS et statistiques temps réel.} Ajouter un
panneau optionnel affichant le FPS instantané, le coût du tick,
et le nombre d'entités vivantes par espèce. Cela faciliterait
le \emph{tuning} et la détection de régressions de performance.

\paragraph{Profils de configuration.} Permettre la sélection
d'un fichier de configuration alternatif via un argument de
ligne de commande \code{--config calme.conf}, et fournir
quelques presets dans \file{presets/} (\file{calme.conf},
\file{chaos.conf}, \file{predator-heavy.conf}).

\paragraph{Sauvegarde et restauration d'état.} Permettre de
\emph{geler} l'état complet de la simulation (positions,
vitesses, états) et de le restaurer ultérieurement, pour
reproduire des scénarios intéressants.

\paragraph{Spatial hashing.} Implémenter une grille uniforme
indexée sur la position des entités, mise à jour à chaque tick.
Les recherches de voisinage deviennent localisées à quelques
cellules.

\paragraph{Rendu Cairo unique.} Remplacer les widgets
\code{GtkPicture} par un seul \code{GtkDrawingArea} dessiné en
Cairo, ce qui éliminerait le coût de \emph{layout} et
permettrait d'atteindre des populations beaucoup plus
importantes.

\paragraph{Mode multijoueur.} Distribuer la simulation sur
plusieurs instances connectées par socket UDP, avec une
synchronisation périodique de l'état. Le déterminisme du pas
fixe rend cet exercice tractable.

\paragraph{Portage Wayland natif et Windows.} GTK4 supporte
nativement Wayland~; un portage Windows via les paquets
\code{MSYS2} est en théorie possible. La principale difficulté
serait le \emph{packaging} et la gestion des chemins de
configuration.

\paragraph{Validation TOML renforcée.} Définir un schéma
explicite des plages acceptables pour chaque paramètre, et
émettre des avertissements quand des valeurs sortent du
raisonnable (par exemple \code{max\_speed > 1000}).

\section{Apport personnel}

Au-delà du résultat livré, le projet a été l'occasion
d'explorer en pratique~:

\begin{itemize}
  \item Le modèle de Reynolds, depuis l'article original de
    1987 jusqu'à ses extensions modernes pour les jeux
    vidéo~\cite{millington2009ai}.
  \item L'écriture d'un parseur de configuration robuste en C,
    avec gestion mémoire stricte et messages d'erreur lisibles.
  \item Les particularités de GTK4 utilisé hors de son contexte
    habituel (interface de bureau classique) pour piloter une
    simulation animée temps réel.
  \item La discipline de \emph{vendoring} d'une bibliothèque
    tierce (\code{tomlc99}) pour garantir la reproductibilité du
    \emph{build} sans dépendance externe.
\end{itemize}

Le \emph{refactoring} qui a fait passer la liste des espèces
d'une \code{enum} compile-time à un registre dynamique
runtime a en particulier été instructif~: il a fallu
identifier précisément tous les points de couplage du code à
l'enum, et les remplacer par des indices entiers résolus à
l'exécution. Cette transformation, en apparence cosmétique,
est en réalité ce qui a rendu possible l'objectif central du
projet (configuration externe).
